\chapter[34]{Constant Mean Curvature Surfaces and Harmonic Maps by IFT}
\label{ch:chow-iv-cmc-surfaces-harmonic-maps-ift}
\dcref{ch34}
\source{chow-et-al-ricci-flow-part-iv:page-0278}

\section{Constant mean curvature hypersurfaces}

Let
\[
 g_{\mathrm{cusp}}=dr^2+e^{-2r}g_{\mathrm{flat}}
\]
on \((a,b)\times\mathcal V^{n-1}\), where \((\mathcal V,g_{\mathrm{flat}})\)
is flat.  If \(\mathcal V_\varphi\) is a graph and the normal points in the
positive \(r\)-direction, the linearization of its mean curvature is
\[
 DH(fN)=\Delta f+(|\mathrm{II}|^2+\Rc(N,N))f.
\]
On a standard cusp slice the zeroth-order terms cancel, so this operator is
the slice Laplacian.

\begin{proposition}[Existence of a CMC sweep-out in almost hyperbolic cusps]
\label{prop:chow-iv-cmc-sweepout-almost-hyperbolic-cusp}
\source{chow-et-al-ricci-flow-part-iv:page-0278,chow-et-al-ricci-flow-part-iv:page-0279,chow-et-al-ricci-flow-part-iv:page-0280}
\dcref{ch34:34.1}
\group{Constant Mean Curvature Hypersurfaces}

Given \([a,b]\subset\mathbb R\) and \(\varepsilon_0>0\), every smooth metric
\(g\) on \([a,b]\times\mathcal V^{n-1}\) which is sufficiently close to
\(g_{\mathrm{cusp}}\) in \(C^{2,\alpha}\), \(0<\alpha<1\), admits a smooth
one-parameter family of smooth \(g\)-CMC hypersurfaces.  They are
\(C^{2,\alpha}\)-close to the standard slices and sweep out
\((a+\varepsilon_0,b-\varepsilon_0)\times\mathcal V\).
\end{proposition}
\begin{proof}
\sketch
For a mean-zero function \(\varphi\) on \(\mathcal V\), set
\[
 \Phi(g,r,\varphi,c)=H_g(\mathcal V_{r+\varphi})+c.
\]
At \((g_{\mathrm{cusp}},r,0,n-1)\), its derivative in \((\varphi,c)\) is
\[
 (f,u)\longmapsto\Delta_r f+u
 :\widehat C^{2,\alpha}(\mathcal V)\times\mathbb R
 \longrightarrow C^\alpha(\mathcal V).
\]
This map is injective because integration first gives \(u=0\), after which a
mean-zero harmonic function is zero.  It is surjective because, for any
\(h\), the equation \(\Delta_r f=h-h_{\mathrm{avg}}\) has a unique mean-zero
solution, with the required Schauder estimate.  The implicit function theorem,
uniformly for \(r\) in a compact subinterval, therefore gives unique smooth
parameters \((\varphi_r,c_r)\) depending smoothly on \(r\) and \(g\).
Elliptic regularity makes \(\varphi_r\) smooth when \(g\) is smooth.  The
mean-zero normalization makes the graphs distinct, while their uniform
closeness to the end slices and continuity in \(r\) show that they cover the
stated smaller product region.
\end{proof}

\begin{remark}[Smooth dependence of a fixed-area CMC slice]
\label{rem:chow-iv-cmc-fixed-area-smooth-dependence}
\source{chow-et-al-ricci-flow-part-iv:page-0280}
\uses{prop:chow-iv-cmc-sweepout-almost-hyperbolic-cusp}
\dcref{ch34:34.2}
\group{Constant Mean Curvature Hypersurfaces}


If \(g(t)\) is a smooth family of metrics close to the cusp metric and a
standard cusp slice has area \(A\), then the corresponding CMC hypersurface of
area \(A\) depends smoothly on \(t\).  Indeed, the area of standard slices has
nonzero derivative in the \(r\)-direction, so the area constraint supplies a
second implicit-function argument.
\end{remark}

\begin{claim}[Foliation by CMC cusp graphs]
\label{clm:chow-iv-cmc-cusp-graphs-foliate}
\source{chow-et-al-ricci-flow-part-iv:page-0280}
\uses{prop:chow-iv-cmc-sweepout-almost-hyperbolic-cusp}
\dcref{ch34:34.4}
\group{Constant Mean Curvature Hypersurfaces}


For \(g\) sufficiently close to \(g_{\mathrm{cusp}}\), a neighborhood of
\([a+2\varepsilon_0,b-2\varepsilon_0]\times\mathcal V\) is foliated by the
CMC graphs \(\mathcal V_{r+\varphi_r}\),
\(r\in(a+\varepsilon_0,b-\varepsilon_0)\).
\end{claim}

\begin{proposition}[Existence of CMC spheres in necks]
\label{prop:chow-iv-cmc-spheres-almost-cylindrical-neck}
\source{chow-et-al-ricci-flow-part-iv:page-0281}
\uses{prop:chow-iv-cmc-sweepout-almost-hyperbolic-cusp}
\dcref{ch34:34.5}
\group{Constant Mean Curvature Hypersurfaces}


Given \([a,b]\subset\mathbb R\) and \(\varepsilon>0\), every smooth metric on
\([a,b]\times S^{n-1}\) sufficiently close in \(C^{2,\alpha}\) to
\(g_{\mathrm{cyl}}=dr^2+g_{\mathrm{sph}}\) admits a smooth family of CMC
hypersurfaces which are close to the standard slices and sweep out
\([a+\varepsilon,b-\varepsilon]\times S^{n-1}\).
\end{proposition}
\begin{proof}
\sketch
For a cylindrical slice, \(\Rc(\partial_r,\partial_r)=0\) and
\(\mathrm{II}=0\).  Thus the derivative of
\(\Phi(g,r,\varphi,c)=H_g(\mathcal S_{r+\varphi})+c\) in the mean-zero graph
and constant variables is again \((f,u)\mapsto\Delta f+u\).  The invertibility,
uniform implicit-function argument, regularity, and sweep-out argument in the
proof of Proposition~\ref{prop:chow-iv-cmc-sweepout-almost-hyperbolic-cusp}
apply verbatim.
\end{proof}

\section{Harmonic self-maps of the sphere}

For a diffeomorphism \(F:(M,g)\to(M,\widetilde g)\), pull its map-Laplacian
back to \(TM\).  At the identity its linearization is
\[
 L_gV=\Delta_gV+\Rc_g(V).
\]

\begin{proposition}[Harmonic maps of the unit sphere near the identity]
\label{prop:chow-iv-harmonic-sphere-map-near-identity}
\source{chow-et-al-ricci-flow-part-iv:page-0281,chow-et-al-ricci-flow-part-iv:page-0285,chow-et-al-ricci-flow-part-iv:page-0286,chow-et-al-ricci-flow-part-iv:page-0287}
\uses{lem:chow-iv-sphere-map-laplacian-kernel, lem:chow-iv-sphere-maps-near-identity-parametrization, lem:chow-iv-map-laplacian-killing-orthogonality}
\dcref{ch34:34.6}
\group{Harmonic Maps on Spheres}


If \(n\ge3\), every harmonic self-map of the unit sphere \(S^n\) sufficiently
close to the identity in \(C^{2,\alpha}\) is an isometry.  If \(n=2\), every
such map is a conformal diffeomorphism.
\end{proposition}
\begin{proof}
\sketch
Parametrize maps near the identity uniquely as \(\Phi(f,\sigma)\), with \(f\)
an isometry and \(\sigma\) orthogonal to Killing fields when \(n\ge3\).  After
identifying the pulled-back map-Laplacian using the two metrics, its value is
orthogonal to Killing fields.  Its derivative in \(\sigma\) at the identity is
\(L_{g_{\mathrm{sph}}}\) restricted to that orthogonal complement.  By the
kernel description and self-adjoint Fredholm theory this restriction is an
isomorphism.  The implicit function theorem says that the unique nearby zero
over every nearby isometry is \(\sigma=0\), proving the first assertion.  For
\(n=2\), replace isometries and Killing fields by conformal diffeomorphisms and
conformal Killing fields; the same argument applies.
\end{proof}

\begin{remark}[Low eigenspaces of the linearized map-Laplacian]
\label{rem:chow-iv-sphere-map-laplacian-low-eigenspaces}
\source{chow-et-al-ricci-flow-part-iv:page-0282,chow-et-al-ricci-flow-part-iv:page-0283}
\dcref{ch34:34.7}
\group{Harmonic Maps on Spheres}

On the unit sphere, the eigenspace of \(L_{g_{\mathrm{sph}}}\) of eigenvalue
\(2-n\) consists of the closed one-forms with Hodge eigenvalue \(n\), its
kernel consists of the coclosed one-forms with Hodge eigenvalue \(2(n-1)\),
and all its other eigenvalues are positive.
\end{remark}

\begin{lemma}[Kernel of the linearized map-Laplacian on the sphere]
\label{lem:chow-iv-sphere-map-laplacian-kernel}
\source{chow-et-al-ricci-flow-part-iv:page-0283,chow-et-al-ricci-flow-part-iv:page-0284}
\uses{rem:chow-iv-sphere-map-laplacian-low-eigenspaces}
\dcref{ch34:34.8}
\group{Harmonic Maps on Spheres}


For the unit sphere, \(\ker L_{g_{\mathrm{sph}}}\) is the space dual to the
Killing vector fields if \(n\ge3\), and is the space dual to the conformal
Killing vector fields if \(n=2\).
\end{lemma}
\begin{proof}
\sketch
For every vector field \(W\) on a closed manifold, integration by parts and
commutation of derivatives give Yano's identity
\[
 -\int_M\langle L_gW,W\rangle
 =\frac12\int_M\left|\mathcal L_Wg-\frac2n(\operatorname{div}W)g\right|^2
  -\left(1-\frac2n\right)\int_M|\operatorname{div}W|^2.
\]
When \(n=2\), this identifies the kernel with conformal Killing fields.  When
\(n\ge3\), the spectral description in Remark~\ref{rem:chow-iv-sphere-map-laplacian-low-eigenspaces}
makes a kernel element coclosed, hence divergence-free, and the identity makes
it Killing.  Conversely, taking the divergence of the Killing equation yields
\(\Delta_gW+\Rc(W)=0\), so every Killing field lies in the kernel.
\end{proof}

\begin{lemma}[Parametrization of maps near the identity]
\label{lem:chow-iv-sphere-maps-near-identity-parametrization}
\source{chow-et-al-ricci-flow-part-iv:page-0284,chow-et-al-ricci-flow-part-iv:page-0285}
\dcref{ch34:34.9}
\group{Harmonic Maps on Spheres}

For \(n\ge2\), a neighborhood of \((\operatorname{id}_{S^n},0)\) in
\[
 \operatorname{Isom}(S^n)\times\operatorname{KV}(S^n)_{2,\alpha}^{\perp}
\]
is mapped diffeomorphically onto a neighborhood of the identity in
\(C^{2,\alpha}(S^n,S^n)\) by
\(\Phi(f,\sigma)(x)=\exp_{f(x)}(\sigma_{f(x)})\).  When \(n=2\), the analogous
statement holds with conformal diffeomorphisms and conformal Killing fields.
\end{lemma}
\begin{proof}
\sketch
At \((\operatorname{id},0)\), the derivative is
\((\tau,\nu)\mapsto\tau+\nu\).  The Killing fields and their \(L^2\)-orthogonal
complement form the whole section space, so this derivative is an isomorphism.
The inverse function theorem proves the first assertion.  The identical
argument with conformal Killing fields proves the second.
\end{proof}

\begin{lemma}[Orthogonality of the map-Laplacian to Killing fields]
\label{lem:chow-iv-map-laplacian-killing-orthogonality}
\source{chow-et-al-ricci-flow-part-iv:page-0285,chow-et-al-ricci-flow-part-iv:page-0286}
\dcref{ch34:34.10}
\group{Harmonic Maps on Spheres}

Let \(F:(M,g)\to(N,h)\) be a diffeomorphism of closed Riemannian manifolds and
let \(I:TM\to TM\) satisfy
\(\langle V,W\rangle_{F^*h}=\langle I(V),W\rangle_g\).  For every Killing
field \(X\) on \((M,g)\),
\[
 \left\langle I((F^{-1})_*\Delta_{g,h}F),X\right\rangle_{L^2(g)}=0.
\]
For \(n=2\), the same is true for every conformal Killing field.
\end{lemma}
\begin{proof}
\sketch
Naturality gives \((F^{-1})_*\Delta_{g,h}F=\Delta_{g,F^*h}\operatorname{id}\).
Writing \(\widetilde g=F^*h\), direct calculation yields
\[
 \Delta_{g,\widetilde g}\operatorname{id}
 =\left(\operatorname{div}_g\widetilde g-	frac12d\operatorname{tr}_g\widetilde g\right)^{\sharp_{\widetilde g}}.
\]
Integration by parts turns the claimed pairing into
\[
 -\frac12\int_M
 \langle\widetilde g,\mathcal L_Xg-(\operatorname{div}X)g\rangle_g.
\]
This vanishes for a Killing field, and in dimension two it also vanishes for a
conformal Killing field.
\end{proof}

\begin{remark}[The two-dimensional harmonic-map argument]
\label{rem:chow-iv-harmonic-two-sphere-near-identity}
\source{chow-et-al-ricci-flow-part-iv:page-0287}
\uses{lem:chow-iv-sphere-map-laplacian-kernel, lem:chow-iv-sphere-maps-near-identity-parametrization, lem:chow-iv-map-laplacian-killing-orthogonality}
\dcref{ch34:34.11}
\group{Harmonic Maps on Spheres}


In dimension two, the proof of Proposition~\ref{prop:chow-iv-harmonic-sphere-map-near-identity}
uses the conformal group, conformal Killing fields, and their orthogonal
complement in place of the corresponding isometric objects.
\end{remark}

\section{Harmonic maps on negatively Ricci-curved manifolds}

\begin{definition}[Normal boundary condition]
\label{def:chow-iv-harmonic-map-normal-boundary-condition}
\source{chow-et-al-ricci-flow-part-iv:page-0287}
\dcref{ch34:34.12}
\group{Harmonic Maps with Boundary}

A harmonic map \(F:(M,g)\to(N,h)\) satisfies the normal boundary condition if
\(F(\partial M)\subset\partial N\) and \(F_*(N_g)\) is normal to \(\partial N\)
with respect to \(h\).
\end{definition}

\begin{proposition}[Existence of harmonic maps near the identity]
\label{prop:chow-iv-harmonic-diffeomorphism-negative-ricci}
\source{chow-et-al-ricci-flow-part-iv:page-0287,chow-et-al-ricci-flow-part-iv:page-0293}
\uses{def:chow-iv-harmonic-map-normal-boundary-condition, lem:chow-iv-linearized-harmonic-boundary-operator-invertible}
\dcref{ch34:34.13}
\group{Harmonic Maps with Boundary}


Let \((M^n,g)\) be compact and smooth, with negative Ricci curvature and
concave boundary.  Every smooth metric \(\widetilde g\) sufficiently close to
\(g\) in \(C^{2,\alpha}\) admits a unique smooth harmonic diffeomorphism
\[
 F:(M,g)\longrightarrow(M,\widetilde g)
\]
which satisfies the normal boundary condition and is \(C^{2,\alpha}\)-close to
the identity.  For a smoothly varying family \(\widetilde g(t)\), the maps
\(F(t)\) vary smoothly as well.
\end{proposition}
\begin{proof}
\sketch
For diffeomorphisms preserving the boundary, form the pulled-back
map-Laplacian together with the tangential component of \(F_*N\).  Its zero set
is precisely the set of harmonic maps satisfying the normal boundary
condition.  The derivative at \((g,\operatorname{id})\) in the map variable is
the isomorphism of Lemma~\ref{lem:chow-iv-linearized-harmonic-boundary-operator-invertible}.
The implicit function theorem gives a unique nearby \(C^{2,\alpha}\) solution
and smooth dependence on parameters.  The harmonic-map system has smooth
coefficients; its boundary symbol is the same mixed Neumann--Dirichlet symbol
used in the regularity lemma.  Interior and boundary elliptic bootstrapping
therefore make \(F\) smooth.
\end{proof}

Let \(W^{1,2}_{\top}(TM)\) consist of Sobolev vector fields whose boundary
trace is tangent.  With inward normal \(N\) and second fundamental form
\(\mathrm{II}\), define
\[
 I(U,V)=\int_M(\langle\nabla U,\nabla V\rangle-\Rc(U,V))\,d\mu
 -\int_{\partial M}\mathrm{II}(U,V)\,d\sigma.
\]
Negative Ricci curvature and boundary concavity make this form coercive.

\begin{lemma}[Weak solutions of the linearized equation]
\label{lem:chow-iv-linearized-harmonic-boundary-weak-solution}
\source{chow-et-al-ricci-flow-part-iv:page-0289,chow-et-al-ricci-flow-part-iv:page-0290}
\dcref{ch34:34.14}
\group{Harmonic Maps with Boundary}

If \((M,g)\) is compact with negative Ricci curvature and concave boundary,
then for every \(Q\in L^2(TM)\) and \(f\in L^2(T\partial M)\) there is a unique
\(U\in W^{1,2}_{\top}(TM)\) such that
\[
 I(U,V)=\int_M\langle Q,V\rangle\,d\mu
       +\int_{\partial M}\langle f,V\rangle\,d\sigma
\]
for every \(V\in W^{1,2}_{\top}(TM)\).  Equivalently, \(U\) weakly solves
\[
 -\Delta U-\Rc(U)=Q,qquad
 (\nabla_NU)_{\top}-\mathrm{II}(U)=f,qquad U_\perp=0.
\]
\end{lemma}
\begin{proof}
\sketch
The trace theorem makes \(I\) bounded.  Since \(-\Rc\) is uniformly positive
and \(-\mathrm{II}\) is nonnegative, \(I(U,U)\ge\delta\|U\|_{W^{1,2}}^2\).
The right-hand side is a bounded functional by Cauchy--Schwarz and the trace
theorem.  Lax--Milgram gives existence and uniqueness; integration by parts
identifies its Euler--Lagrange equation and boundary condition.
\end{proof}

\begin{lemma}[Regularity for the linearized mixed boundary problem]
\label{lem:chow-iv-linearized-harmonic-boundary-regularity}
\source{chow-et-al-ricci-flow-part-iv:page-0290,chow-et-al-ricci-flow-part-iv:page-0291}
\uses{lem:chow-iv-linearized-harmonic-boundary-weak-solution}
\dcref{ch34:34.15}
\group{Harmonic Maps with Boundary}


Under the hypotheses of Lemma~\ref{lem:chow-iv-linearized-harmonic-boundary-weak-solution},
if \(Q\) and \(f\) are smooth, then the unique weak solution is smooth.
\end{lemma}
\begin{proof}
\sketch
In boundary normal coordinates, the principal interior symbol is
\[
 A(\xi)\operatorname{id}_{TM},
\]
while the principal boundary symbol is
\[
 V\longmapsto \xi(N)V_{\top}+V_\perp.
\]
For every nonzero boundary covector \(\zeta\), the characteristic polynomial
has conjugate roots \(\tau^\pm(\zeta)\), each of multiplicity \(n\).  Reducing
the rows of the composed boundary symbol and adjugate interior symbol modulo
\((\tau-\tau^+)^n\) leaves the independent tangential rows
\((\tau-\tau^+)^{n-1}\tau\) and the independent normal row
\((\tau-\tau^+)^{n-1}\).  Hence the complementing boundary condition holds.
Regularity for elliptic systems with complementing boundary conditions now
implies smoothness.
\end{proof}

\begin{lemma}[Schauder estimate]
\label{lem:chow-iv-linearized-harmonic-boundary-schauder}
\source{chow-et-al-ricci-flow-part-iv:page-0291,chow-et-al-ricci-flow-part-iv:page-0292}
\dcref{ch34:34.16}
\group{Harmonic Maps with Boundary}

For every \(U\in C^{2,\alpha}(TM)\) with \(U_\perp=0\) on \(\partial M\),
\[
 \|U\|_{C^{2,\alpha}}
 \le C\|{-\Delta U-\Rc(U)}\|_{C^\alpha}
 +C\|(\nabla_NU)_{\top}-\mathrm{II}(U)\|_{C^{1,\alpha}(T\partial M)}.
\]
\end{lemma}
\begin{proof}
\sketch
The complementing boundary condition gives the usual estimate with an
additional \(C^0\)-term.  Interpolation bounds the \(C^2\)-norm by a small
multiple of the top Hölder seminorm plus the \(C^0\)-norm, and compactness gives
\(\|U\|_{C^0}\le\varepsilon\|\nabla U\|_{C^0}+C_\varepsilon\|U\|_{L^2}\).
Finally, coercivity of \(I\), Cauchy--Schwarz, and the trace theorem bound the
\(L^2\)-term by the interior and boundary data.  Choosing the small constants
so that derivative terms are absorbed proves the estimate.
\end{proof}

\begin{lemma}[Invertibility of the linearized harmonic boundary operator]
\label{lem:chow-iv-linearized-harmonic-boundary-operator-invertible}
\source{chow-et-al-ricci-flow-part-iv:page-0293}
\uses{lem:chow-iv-linearized-harmonic-boundary-weak-solution, lem:chow-iv-linearized-harmonic-boundary-regularity, lem:chow-iv-linearized-harmonic-boundary-schauder}
\dcref{ch34:34.17}
\group{Harmonic Maps with Boundary}


The operator
\[
 U\longmapsto
 \bigl(\Delta U+\Rc(U),(\nabla_NU)_{\top}-\mathrm{II}(U)\bigr)
\]
from tangent-boundary \(C^{2,\alpha}\) vector fields to
\(C^\alpha(TM)\times C^{1,\alpha}(T\partial M)\) is an isomorphism.
\end{lemma}
\begin{proof}
\sketch
For smooth data, weak existence and regularity give a smooth solution.  Approximate
arbitrary Hölder data by smooth data.  The Schauder estimate makes the associated
solutions Cauchy in \(C^{2,\alpha}\), so their limit solves the prescribed
problem.  The same estimate gives uniqueness and boundedness of the inverse.
\end{proof}

\begin{claim}[Harmonic diffeomorphism near an approximate isometry]
\label{clm:chow-iv-harmonic-diffeomorphism-near-approximate-isometry}
\source{chow-et-al-ricci-flow-part-iv:page-0293}
\uses{prop:chow-iv-harmonic-diffeomorphism-negative-ricci}
\dcref{ch34:34.18}
\group{Harmonic Maps with Boundary}


An approximate isometry between compact manifolds satisfying the negative
Ricci-curvature and concave-boundary hypotheses can, when sufficiently close,
be followed by a harmonic diffeomorphism close to it and satisfying the normal
boundary condition.
\end{claim}

\section{Almost isometries of hyperbolic manifolds}

\begin{lemma}[Dependence of hypersurface geometry on the ambient metric]
\label{lem:chow-iv-hypersurface-geometry-ambient-metric-stability}
\source{chow-et-al-ricci-flow-part-iv:page-0294,chow-et-al-ricci-flow-part-iv:page-0295,chow-et-al-ricci-flow-part-iv:page-0296}
\dcref{ch34:34.19}
\group{Hyperbolic Almost Isometries}

Let \(N^{n-1}\subset M^n\) be an oriented smooth hypersurface, and let \(\nu\)
and \(\widetilde\nu\) be the same-side unit normals for metrics \(g\) and
\(\widetilde g\).  Uniformly on a fixed coordinate neighborhood,
\begin{align*}
 |\widetilde\nu-\nu|_g&\le C\|\widetilde g-g\|_{C^0},\\
 |\widetilde\nabla\widetilde\nu|_g&\le C,
 &|\nabla(\widetilde\nu-\nu)|_g&\le C\|\widetilde g-g\|_{C^1},\\
 |\widetilde{\mathrm{II}}(X,Y)-\mathrm{II}(X,Y)|
 &\le C|X|_g|Y|_g\|\widetilde g-g\|_{C^1},\\
 |\widetilde\nabla\widetilde{\mathrm{II}}(X,Y)-\nabla\mathrm{II}(X,Y)|
 &\le C|X|_g|Y|_g\|\widetilde g-g\|_{C^2}.
\end{align*}
\end{lemma}
\begin{proof}
\sketch
Choose adapted coordinates with \(N=\{x^n=0\}\).  Solving the orthogonality
equations expresses the unnormalized normal rationally in the metric
coefficients and the inverse tangential metric; normalization then proves the
first estimate, and differentiation proves the second pair.  The identity
\(\mathrm{II}(X,Y)=g(\nabla_X\nu,Y)\), together with
\(|(\widetilde\nabla-\nabla)Y|_g\le C|Y|_g\|\widetilde g-g\|_{C^1}\), gives
the third estimate.  Differentiating this identity once more gives the last.
\end{proof}

\begin{lemma}[Stability of intrinsic hypersurface curvature]
\label{lem:chow-iv-hypersurface-intrinsic-curvature-stability}
\source{chow-et-al-ricci-flow-part-iv:page-0296}
\uses{lem:chow-iv-hypersurface-geometry-ambient-metric-stability}
\dcref{ch34:34.20}
\group{Hyperbolic Almost Isometries}


With the hypotheses above, if \(g_N\) and \(\widetilde g_N\) are the two induced
metrics, then
\[
 |\Rm_{g_N}-\Rm_{\widetilde g_N}|
 \le C\|\widetilde g-g\|_{C^2(M,g)}.
\]
\end{lemma}
\begin{proof}
\sketch
Use the Gauss formula
\(\nabla^N_XY=\nabla_XY+\mathrm{II}(X,Y)\nu\) for both metrics.  Expanding the
two intrinsic curvature tensors in terms of these connections and applying
Lemma~\ref{lem:chow-iv-hypersurface-geometry-ambient-metric-stability} gives
the estimate.
\end{proof}

\begin{remark}[Higher-order ambient metric stability]
\label{rem:chow-iv-hypersurface-geometry-higher-order-stability}
\source{chow-et-al-ricci-flow-part-iv:page-0296}
\uses{lem:chow-iv-hypersurface-geometry-ambient-metric-stability, lem:chow-iv-hypersurface-intrinsic-curvature-stability}
\dcref{ch34:34.21}
\group{Hyperbolic Almost Isometries}


For \(k\ge1\), the corresponding \(k\)-th derivatives of the second
fundamental forms are controlled by \(C_k\|\widetilde g-g\|_{C^{k+1}}\), and
the \(k\)-th derivatives of the intrinsic curvature tensors are controlled by
\(C_k\|\widetilde g-g\|_{C^{k+2}}\).  The same statements hold when the
hypersurface is a boundary component.
\end{remark}

\begin{theorem}[Existence of isometries near almost isometries]
\label{thm:chow-iv-isometry-near-almost-isometry}
\source{chow-et-al-ricci-flow-part-iv:page-0296,chow-et-al-ricci-flow-part-iv:page-0297,chow-et-al-ricci-flow-part-iv:page-0298,chow-et-al-ricci-flow-part-iv:page-0299}
\uses{lem:chow-iv-hypersurface-geometry-ambient-metric-stability, lem:chow-iv-hypersurface-intrinsic-curvature-stability}
\dcref{ch34:34.22}
\group{Hyperbolic Almost Isometries}


Let \((\mathcal H^3,h)\) be a finite-volume hyperbolic manifold and
\(\mathcal H_A\) a sufficiently deep standard cusp truncation.  For every
\(\ell\in\mathbb N\) there is \(k=k(\ell)\) such that, whenever
\((\widetilde{\mathcal H}^3,\widetilde h)\) is finite-volume hyperbolic with at
least as many cusp ends and an embedding
\(F:\mathcal H_A\to\widetilde{\mathcal H}\) satisfies
\[
 \|F^*\widetilde h-h\|_{C^k(\mathcal H_A,h)}<\frac1k,
\]
there is an isometry \(I:(\mathcal H,h)\to
(\widetilde{\mathcal H},\widetilde h)\) with
\[
 d_{C^\ell(\mathcal H_A,h)}(F,I)<\frac1\ell.
\]
In particular, the two complete hyperbolic manifolds are isometric.
\end{theorem}
\begin{proof}
\sketch
First fill every image boundary torus which is not a standard truncating torus
by the compact solid torus or ball-side component it bounds.  The resulting
compact submanifold has boundary in the boundary of the target truncation.  If
the target had an extra cusp, or if one image torus had been filled, an
untruncated end would remain, contradicting compactness.  Thus the two
manifolds are diffeomorphic, and Mostow rigidity makes them isometric; identify
them henceforth.

Every isometric embedding of \(\mathcal H_A\) into \(\mathcal H\) extends to a
global isometry.  Indeed, its boundary tori lie in the Margulis thin part.
An isometry carries the intrinsically flat, totally umbilic foliation of each
truncated cusp to the corresponding standard foliation.  In cusp coordinates
it is therefore a product of a linear radial map and a flat-torus homothety,
which extends uniquely along the whole cusp.

If the asserted quantitative conclusion failed, there would be increasingly
accurate almost-isometric embeddings \(F_i\) staying a fixed \(C^\ell\)-distance
from every isometry.  The two stability lemmas show that their image boundary
tori become intrinsically flat and totally umbilic, hence lie arbitrarily close
to the fixed truncation boundary.  Compactness and the estimates then yield a
subsequence converging smoothly to a local isometry
\(F_\infty:\mathcal H_A\to\mathcal H_A\).  The inverse function theorem and
injectivity of the \(F_i\) make the interior limit injective, and boundary
convergence makes its boundary restriction a diffeomorphism.  Thus
\(F_\infty\) is an isometric embedding, so the preceding paragraph extends it
to a global isometry \(I_\infty\).  Smooth convergence contradicts the fixed
positive distance from all isometries.
\end{proof}

\begin{remark}[Foliation question]
\label{rem:chow-iv-cmc-sweepout-foliation-question}
\dcref{ch34:34.3}
\group{Constant Mean Curvature Hypersurfaces}
\source{chow-et-al-ricci-flow-part-iv:page-0280}
\uses{prop:chow-iv-cmc-sweepout-almost-hyperbolic-cusp}
The sweep-out should in fact be a foliation; this stronger conclusion is not
needed for the Ricci-flow applications.
\end{remark}
